% Einheit 2 von 3 – Zehnerpotenzen (bank/potenzen-wurzeln/e2.jsonl, zusammenbau v0.1)
%% TODO Zweigzeile Teil 1: was hier gelernt wird (Satz in Schülersprache aus dem Einheitstitel) – Einheit 2
\einheitenkopf[e2]{Einheit 2 von 3 · Zehnerpotenzen}
\zweigzeile{ab Kl. 9, je nach Buch bis Kl. 10 (am Gymnasium ab Kl. 8) · P10}
\setcounter{aufgabe}{18}

%% TODO Titel: Ich-kann-Satz fehlt in der Bank (Platzhalter: Kettenname) – Nr. 19
%% TODO Anweisung: ein Satz über der ersten Teilaufgabe fehlt in der Bank – Nr. 19
\begin{aufgabe}{Zehnerpotenzen}
\begin{teile}
\teil $4{,}7 \cdot 10^{8}$ – größer oder kleiner als eins?\\ \kreuz{größer als eins} \\ \kreuz{kleiner als eins}
\end{teile}
%% TODO Darstellung für schwach fehlt in der Bank (potenzen-wurzeln-e2-k1-s1-v1; Abschnitt „Für schwache Schüler“)
\swz{$10^{4}$ – ausgeschrieben?}{}
%% TODO Darstellung für schwach fehlt in der Bank (potenzen-wurzeln-e2-k1-s1-v2; Abschnitt „Für schwache Schüler“)
\swz{$10^{7}$ – ausgeschrieben?}{}
%% TODO Darstellung für schwach fehlt in der Bank (potenzen-wurzeln-e2-k1-s1-v3; Abschnitt „Für schwache Schüler“)
\swz{$1\,000$ – als Zehnerpotenz?}{}
%% TODO Darstellung für schwach fehlt in der Bank (potenzen-wurzeln-e2-k1-s1-v4; Abschnitt „Für schwache Schüler“)
\swz{$1\,000\,000\,000$ – als Zehnerpotenz?}{}
%% TODO Darstellung für schwach fehlt in der Bank (potenzen-wurzeln-e2-k1-s1-v5; Abschnitt „Für schwache Schüler“)
\swz{$100$ – als Zehnerpotenz?}{}
\swfrage{Was bleibt gleich, was ändert sich?}
\begin{teile}
\teil $10^{9}$ – welches Zahlwort?\\ \kreuz{Million} \\ \kreuz{Milliarde} \\ \kreuz{Billion}
\end{teile}
%% TODO Darstellung für schwach fehlt in der Bank (potenzen-wurzeln-e2-k1-s3-v1; Abschnitt „Für schwache Schüler“)
\swz{$53 \cdot 100$ – wie viel?}{}
%% TODO Darstellung für schwach fehlt in der Bank (potenzen-wurzeln-e2-k1-s4-v1; Abschnitt „Für schwache Schüler“)
\swz{$29 \cdot 10^{5}$ – wie viel?}{}
%% TODO Darstellung für schwach fehlt in der Bank (potenzen-wurzeln-e2-k1-s5-v1; Abschnitt „Für schwache Schüler“)
\swz{$4{,}8 \cdot 10^{3}$ – wie viel?}{}
%% TODO Darstellung für schwach fehlt in der Bank (potenzen-wurzeln-e2-k1-s6-v1; Abschnitt „Für schwache Schüler“)
\swz{$6{,}3 \cdot 10^{5}$ – ausgeschrieben? \hfill (P10 2016 OS)}{}
%% TODO Darstellung für schwach fehlt in der Bank (potenzen-wurzeln-e2-k1-s7-v1; Abschnitt „Für schwache Schüler“)
\swz{$58\,000$ – in Zehnerpotenzschreibweise?}{}
%% TODO Darstellung für schwach fehlt in der Bank (potenzen-wurzeln-e2-k1-s8-v1; Abschnitt „Für schwache Schüler“)
\swz{$64\,000 = 6{,}4 \cdot 10^{\square}$ – welche Hochzahl?}{}
\begin{teile}
\teil Das Gehirn eines Menschen hat etwa $8{,}6 \cdot 10^{10}$ Nervenzellen. Welche Zahl passt? Kreuze an. \hfill (P10 2015 OS)\\ \kreuz{$860\,000\,000\,000$} \\ \kreuz{$86\,000\,000\,000$} \\ \kreuz{$8\,600\,000\,000$}
\end{teile}
%% TODO Darstellung für schwach fehlt in der Bank (potenzen-wurzeln-e2-k1-s10-v1; Abschnitt „Für schwache Schüler“)
\swz{$3{,}6 \cdot 10^{-3}$ – als Dezimalzahl?}{}
%% TODO Darstellung für schwach fehlt in der Bank (potenzen-wurzeln-e2-k1-s11-v1; Abschnitt „Für schwache Schüler“)
\swz{$0{,}0052$ – in Zehnerpotenzschreibweise?}{}
%% TODO Darstellung für schwach fehlt in der Bank (potenzen-wurzeln-e2-k1-s12-v1; Abschnitt „Für schwache Schüler“)
\swz{$4{,}1 \cdot 10^{5}$; $9{,}3 \cdot 10^{4}$; $2{,}8 \cdot 10^{6}$ – aufsteigend geordnet?}{}
%% TODO Darstellung für schwach fehlt in der Bank (potenzen-wurzeln-e2-k1-s13-v1; Abschnitt „Für schwache Schüler“)
\swz{Der Taschenrechner zeigt 6.2E7 – welche Zahl ist das?}{}
%% TODO Darstellung für schwach fehlt in der Bank (potenzen-wurzeln-e2-k1-s14-v1; Abschnitt „Für schwache Schüler“)
\swz{Von der Erde zur Sonne sind es etwa $149\,597\,871$ km – in Zehnerpotenzschreibweise, Faktor auf zwei Stellen nach dem Komma gerundet?}{}
%% TODO Darstellung für schwach fehlt in der Bank (potenzen-wurzeln-e2-k1-s15-v1; Abschnitt „Für schwache Schüler“)
\swz[3]{Ein Bakterium ist etwa $2 \cdot 10^{-3}$ mm lang. Wie viele liegen hintereinander auf $1$ mm?}{}
%% TODO Darstellung für schwach fehlt in der Bank (potenzen-wurzeln-e2-k1-s16-v1; Abschnitt „Für schwache Schüler“)
\swz{$10^{4} \cdot 10^{3}$ – als Zehnerpotenz?}{}
\swz[3]{$730\,000 = 7{,}3 \cdot 10^{\square}$ – welche Hochzahl? \hfill (P10 2025 OS)}{}
\end{aufgabe}

%% TODO Titel: Ich-kann-Satz fehlt in der Bank (Platzhalter: Kettenname) – Nr. 20
%% TODO Anweisung: ein Satz über der ersten Teilaufgabe fehlt in der Bank – Nr. 20
\begin{aufgabe}{Zehnerpotenzen – Fehler finden, Begründen}
\swz[3]{Ben schreibt $920\,000 = 9{,}2 \cdot 10^{4}$. Wo ist der Fehler?}{}
\swfrage{Begründe, warum $2{,}6 \cdot 10^{5}$ und $26 \cdot 10^{4}$ dieselbe Zahl sind.}
\end{aufgabe}

\uebersichtskasten{Zehnerpotenz: 10ⁿ ist eine Eins mit n Nullen – die Hochzahl zählt die Nullen: 10⁶ = 1 000 000 (eine Million). \\ Mal Zehnerpotenz: Beim Malnehmen mit 10ⁿ hängst du an eine ganze Zahl n Nullen an; bei einer Dezimalzahl wandert das Komma n Stellen nach rechts: 47 · 10⁴ = 470 000; 3,25 · 10⁴ = 32 500. \\ Zehnerpotenzschreibweise: Eine Zahl zwischen eins und zehn mal eine Zehnerpotenz – die Hochzahl zählt, um wie viele Stellen das Komma wandert, nicht die Nullen: 6 700 000 = 6,7 · 10⁶ (das Komma wandert sechs Stellen, Nullen sind es fünf). \\ Negative Hochzahl: 10⁻ⁿ macht die Zahl klein – das Komma wandert n Stellen nach links, vorn kommen Nullen: 4,2 · 10⁻³ = 0,0042.}
